\chapter{Basic Configuration --- Selected ISOs}

\section{Selected ISOs}

This page documents the ISO images selected for the TFG lab environment and the rationale behind each choice.

\begin{center}\rule{0.5\linewidth}{0.5pt}\end{center}

\subsection{ISO overview}

{\footnotesize\begin{longtable}[]{@{}
  >{\raggedright\arraybackslash}p{0.22\linewidth}
  >{\raggedright\arraybackslash}p{0.22\linewidth}
  >{\raggedright\arraybackslash}p{0.22\linewidth}
  >{\raggedright\arraybackslash}p{0.22\linewidth}@{}}
\toprule\noalign{}
\begin{minipage}[b]{\linewidth}\raggedright
Image
\end{minipage} & \begin{minipage}[b]{\linewidth}\raggedright
Version
\end{minipage} & \begin{minipage}[b]{\linewidth}\raggedright
Source
\end{minipage} & \begin{minipage}[b]{\linewidth}\raggedright
Role in labs
\end{minipage} \\
\midrule\noalign{}
\endhead
\bottomrule\noalign{}
\endlastfoot
\textbf{VyOS} & rolling & \href{https://vyos.io}{vyos.io} & Core router --- routing, NAT, firewall \\
\textbf{pfSense CE} & 2.7.x & \href{https://www.pfsense.org}{pfsense.org} & Perimeter firewall, DNS resolver \\
\textbf{Ubuntu Server} & 24.04 LTS & \href{https://ubuntu.com}{ubuntu.com} & Target server --- nginx, SSH \\
\textbf{Ubuntu Desktop} & 24.04 LTS & \href{https://ubuntu.com}{ubuntu.com} & Victim workstation \\
\textbf{Parrot OS} & rolling & \href{https://parrotsec.org}{parrotsec.org} & Attacker node \\
\end{longtable}}

\begin{center}\rule{0.5\linewidth}{0.5pt}\end{center}

\subsection{Selection criteria}

All ISOs were chosen based on three criteria:

\begin{enumerate}
\def\labelenumi{\arabic{enumi}.}
\tightlist
\item
  \textbf{Open source / freely redistributable} --- no licensing cost for students or the institution.
\item
  \textbf{EVE-NG compatibility} --- tested boot under QEMU with virtio drivers; added with the \texttt{linux-} prefix under \texttt{/opt/unetlab/addons/qemu/} as required by EVE-NG.
\item
  \textbf{Industry relevance} --- tools and OS families students will encounter in real engagements (pfSense/VyOS in SMB network gear, Ubuntu Server as the dominant Linux server target, Parrot/Kali as standard pentesting platforms).
\end{enumerate}

\begin{center}\rule{0.5\linewidth}{0.5pt}\end{center}

\subsection{Import into EVE-NG}

\begin{lstlisting}
# 1. Upload ISO to EVE-NG host
scp <image>.iso root@<eveng-host>:/tmp/

# 2. Create the image directory (use linux- prefix)
ssh root@<eveng-host>
mkdir -p /opt/unetlab/addons/qemu/linux-<name>/
mv /tmp/<image>.iso /opt/unetlab/addons/qemu/linux-<name>/cdrom.iso

# 3. Create a virtual disk
qemu-img create -f qcow2 \
  /opt/unetlab/addons/qemu/linux-<name>/virtioa.qcow2 <size>G

# 4. Fix permissions
/opt/unetlab/wrappers/unl_wrapper -a fixpermissions
\end{lstlisting}

\begin{lstlisting}
The `src/scripts/automation/Iso_uploader.py` script in the repository automates
all four steps above with a GUI file picker and automatic disk-size suggestions
based on the OS type.
\end{lstlisting}

\begin{center}\rule{0.5\linewidth}{0.5pt}\end{center}

\subsection{Notes per image}

\begin{description}
\tightlist
\item[\textbf{VyOS (rolling)}]
Must be installed to disk from the live ISO before use in labs. Run \texttt{install\ image} from the VyOS shell after first boot.
\item[\textbf{pfSense CE}]
On first boot, assign interfaces manually via the console menu (option 1). \texttt{vtnet0} $\rightarrow$ LAN, \texttt{vtnet1} $\rightarrow$ WAN in EVE-NG.
\item[\textbf{Ubuntu Server / Desktop}]
Use the \texttt{linux-ubuntu-*} prefix. virtio disk and network drivers are supported out of the box.
\item[\textbf{Parrot OS}]
Use the Security edition. The Home edition does not include pentesting tools.
\end{description}
